\documentclass[11pt]{article}
\usepackage[margin=1in]{geometry}
\usepackage{amsmath}
\usepackage{graphicx}
\usepackage{booktabs}
\usepackage{hyperref}
\usepackage{float}

\begin{document}

\noindent
\textbf{March 6, 2026} \hfill \textbf{Joseph DeLeonardis} \\
\textbf{UIN: 0820000866} \\
\textbf{CSCE-636-700} \\
\textbf{Deep Learning Homework \#5}

\section*{Problem 1}

\subsection*{Chapter 10: Interpreting what ConvNets learn}
This chapter built upon the model created in Chapter 8 using the dogs and cats dataset. Running the Chapter 10 coding exercises demonstrated three key concepts in ConvNet interpretability: layers, gradients, and heatmaps. As a network goes deeper, learned features become increasingly abstract early layers capture basic edges and textures while deeper layers respond to complex concepts like facial features. Grad-CAM heatmaps tie this together by revealing which regions of an image most influenced the model's prediction, providing a visual sanity check that the model is making decisions for the right reasons.
\newline
Sections 1-3 from the homework guidelines are shown in the .ipynb file included. 

\subsection*{Chapter 11: Image Segmentation}
Chapter 11 covers the three types of image segmentation and the distinctions between them. To satisfy the requirements of problem 1 part 4, I modified the image selection after loading the oxford\_segmentation model by replacing the hardcoded index with a randomly generated one, ensuring the segmentation results are drawn from a random sample rather than a fixed image.


\section*{Problem 2}


\subsection*{Chapter 13: }



\end{document}
